% This file is generated from the corresponding Obsidian note.
% Source: Teoricos/GS - Teo12.md
\chapter{Rango constante y rebanadas}
\label{ch:teo-12}
\section{Teorema del rango constante}
\begin{bookdefinition}{Funcion de rango constante}
Una funcion suave

\[
F:M\to N
\]
se dice de rango constante $r$ si para todo $p\in M$, la imagen de

\[
(dF)_p
\]
tiene dimension $r$.

\end{bookdefinition}
\begin{bookdefinition}{Submersion}
Una funcion suave

\[
F:M\to N
\]
se dice submersion suave si

\[
(dF)_p
\]
es sobreyectiva para todo $p\in M$.

\end{bookdefinition}
\begin{bookexample}{Submersion}
Sean $m\ge n$. Considerar

\[
\pi:\mathbb R^m\to \mathbb R^n,
\]
\[
(s_1,\ldots,s_n,s_{n+1},\ldots,s_m)\mapsto (s_1,\ldots,s_n).
\]
Notar que $\pi$ es transformacion lineal.

\end{bookexample}
\begin{bookexample}{Rango constante}
\begin{enumerate}
\item Si
\end{enumerate}
\[
F:M\to N
\]
es una inmersion, entonces $\ker(dF)_p$ es trivial para todo $p\in M$, entonces $F$ es de rango constante e igual a $\dim M$. 2. Si

\[
F:M\to N
\]
es submersion, entonces $F$ es de rango constante e igual a $\dim N$. 3. Sean $m,n\in\mathbb N$, $m,n\ge 2$, y sea $r<\min\{m,n\}$. La transformacion lineal

\[
T:\mathbb R^m\to \mathbb R^n,
\]
\[
(s_1,\ldots,s_r,s_{r+1},\ldots,s_m)\mapsto (s_1,\ldots,s_r,0,\ldots,0)
\]
es una funcion de rango constante $r$. Notar que $T$ es transformacion lineal.

\end{bookexample}
\begin{booktheorem}{Forma local de una funcion de rango constante}
Sean $M$ y $N$ variedades suaves de dimension $m$ y $n$ respectivamente, y sea

\[
F:M\to N
\]
una funcion suave de rango constante $r$. Entonces, para cada $p\in M$, existen cartas suaves $(U,\varphi)$ de $M$ centrada en $p$ y $(V,\psi)$ de $N$ centrada en $F(p)$ tales que

\[
F(U)\subseteq V
\]
y

\[
\psi\circ F\circ \varphi^{-1}(s_1,\ldots,s_r,s_{r+1},\ldots,s_m)=(s_1,\ldots,s_r,0,\ldots,0).
\]
En particular, si $F$ es una submersion suave, entonces $(r=n)$ y

\[
\psi\circ F\circ \varphi^{-1}(s_1,\ldots,s_m)=(s_1,\ldots,s_n),
\]
y si $F$ es una inmersion suave, entonces $(r=m)$ y

\[
\psi\circ F\circ \varphi^{-1}(s_1,\ldots,s_m)=(s_1,\ldots,s_m,0,\ldots,0).
\]
\end{booktheorem}
\begin{bookproof}{Revisar}
\begin{enumerate}
\item Sea $p\in M$. Tomamos cartas suaves $(U,\varphi)$ de $M$ centrada en $p$ y $(V,\psi)$ de $N$ centrada en $F(p)$ tales que $F(U)\subseteq V$.
\item Consideramos
\end{enumerate}
\[
\widetilde F=\psi\circ F\circ\varphi^{-1}:\varphi(U)\subseteq\mathbb R^m\to\psi(V)\subseteq\mathbb R^n.
\]
\begin{enumerate}[start=3]
\item Como $\widetilde F$ tiene rango $r$, achicando si hace falta, podemos suponer que una submatriz $r\times r$ de la matriz jacobiana es no singular en $(0,0)$. Reordenando coordenadas, tomamos
\end{enumerate}
\[
\left(\frac{\partial \widetilde F_i}{\partial x_j}(0,0)\right)_{1\le i,j\le r}.
\]
\begin{enumerate}[start=4]
\item Escribimos coordenadas en $\mathbb R^m$ como
\end{enumerate}
\[
s=(x,y),\qquad x=(x_1,\ldots,x_r),\qquad y=(y_1,\ldots,y_{m-r}).
\]
\begin{enumerate}[start=5]
\item Definimos
\end{enumerate}
\[
Q:\widetilde U\subseteq\mathbb R^m\to\mathbb R^m,
\]
\[
Q(x,y)=\big(\widetilde F_1(x,y),\ldots,\widetilde F_r(x,y),y\big).
\]
\begin{enumerate}[start=6]
\item Afirmacion:
\end{enumerate}
\[
dQ(0,0)=\begin{pmatrix}\dfrac{\partial Q_i}{\partial x_j} & \dfrac{\partial Q_i}{\partial y_j}\\0 & I_{m-r}\end{pmatrix}_{(0,0)}
\]
es no singular. Como el bloque

\[
\left(\frac{\partial Q_i}{\partial x_j}(0,0)\right)
\]
es no singular, $dQ(0,0)$ es no singular. 7. Por el teorema de la funcion inversa, existen abiertos $\widetilde U_0$ de $\mathbb R^m$ en $(0,0)$ y $\widehat U$ de $(0,0)$ tales que

\[
Q|_{\widetilde U_0}:\widetilde U_0\to \widehat U
\]
es un difeomorfismo. 8. Como las cajas abiertas son una base de la topologia de $\mathbb R^m$, podemos suponer que $\widehat U$ es una caja:

\[
C_\varepsilon^m(0,0)=\{(s_1,\ldots,s_m)\in\mathbb R^m:\ |s_i|<\varepsilon\}.
\]
\begin{enumerate}[start=9]
\item Sea
\end{enumerate}
\[
\varphi^{-1}:C_\varepsilon^m(0,0)\subseteq\mathbb R^m\to \widetilde U_0\subseteq\mathbb R^m,
\]
\[
(x,y)\mapsto (A(x,y),B(x,y)).
\]
\begin{enumerate}[start=10]
\item Por definicion,
\end{enumerate}
\[
Q(A(x,y),B(x,y))=(x,y).
\]
\begin{enumerate}[start=11]
\item Esto dice que
\end{enumerate}
\[
\big(\widetilde F_1(A,B),\ldots,\widetilde F_r(A,B),B(x,y)\big)=(x,y),
\]
\begin{enumerate}[start=12]
\item Luego $B(x,y)=y$ y entonces
\end{enumerate}
\[
\varphi^{-1}(x,y)=(A(x,y),y).
\]
\begin{enumerate}[start=13]
\item Con esto
\end{enumerate}
\[
\widetilde F\circ\varphi^{-1}(x,y)=\widetilde F(A(x,y),y)=(x,\widetilde B(x,y)),
\]
donde

\[
\widetilde B:C_\varepsilon^m(0,0)\subseteq\mathbb R^m\to\mathbb R^{n-r}.
\]
\begin{enumerate}[start=14]
\item Como $\widetilde F\circ\varphi^{-1}$ tiene rango $r$, y la primera parte ya aporta las $r$ coordenadas $x$, se tiene
\end{enumerate}
\[
\frac{\partial \widetilde B_j}{\partial y_i}(x,y)=0
\]
para todo $(x,y)\in C_\varepsilon^m(0,0)$. 15. Es decir, $\widetilde B_j$ no depende de las variables $y_1,\ldots,y_{m-r}$. En particular, como $C_\varepsilon^m(0,0)$ es una caja,

\[
\widetilde B_j(x,y)=\widetilde B_j(x,0).
\]
\begin{enumerate}[start=16]
\item Definimos
\end{enumerate}
\[
S:C_\varepsilon^r(0)\subseteq\mathbb R^r\to\mathbb R^{n-r},\qquad S(x)=\widetilde B(x,0).
\]
\begin{enumerate}[start=17]
\item Entonces
\end{enumerate}
\[
\widetilde F\circ\varphi^{-1}(x,y)=(x,S(x)).
\]
\begin{enumerate}[start=18]
\item Ahora definimos una carta en el codominio:
\end{enumerate}
\[
\psi(w,u)=(w,u-S(w)).
\]
\begin{enumerate}[start=19]
\item Esta funcion es un difeomorfismo local, pues
\end{enumerate}
\[
d\psi(0)=\begin{pmatrix} I_r & 0\\-\dfrac{\partial S_i}{\partial w_j}(0) & I_{n-r}\end{pmatrix}
\]
es no singular. 20. Finalmente,

\[
\psi\circ\widetilde F\circ\varphi^{-1}(x,y)=\psi(x,S(x))=(x,0).
\]
\bookqed
\end{bookproof}
\begin{bookproposition}{}
Si $X$ es un espacio topologico N2, entonces todo cubrimiento abierto de $X$ tiene un subcubrimiento numerable.

\end{bookproposition}
\begin{booktheorem}{Teorema de Baire}
Si $X$ es un espacio topologico localmente compacto y $T_2$, entonces $X$ no se puede expresar como una union numerable de cerrados con interior vacio.

\end{booktheorem}
\begin{booktheorem}{Del rango global}
Sean $M$ y $N$ variedades suaves, y sea

\[
F:M\to N
\]
funcion suave de rango constante. Entonces:

\begin{itemize}
\item (a)  Si $F$ es sobreyectiva, entonces $F$ es submersion.
\item (b)  Si $F$ es inyectiva, entonces $F$ es inmersion.
\item (c)  Si $F$ es biyectiva, entonces $F$ es difeomorfismo.
\end{itemize}
\end{booktheorem}
\begin{bookproof}{Revisar}
\begin{itemize}
\item (a)
\end{itemize}
\begin{enumerate}
\item Razonemos por absurdo. Si $F$ es sobreyectiva y no es submersion, entonces el rango constante $r$ satisface $r<n$.
\item Por el teorema de rango constante, para cada $p\in M$ existen cartas $(U_p,\varphi_p)$ y $(V_p,\psi_p)$ tales que
\end{enumerate}
\[
\psi_p\circ F\circ\varphi_p^{-1}(x_1,\ldots,x_r,x_{r+1},\ldots,x_m)=(x_1,\ldots,x_r,0,\ldots,0).
\]
\begin{enumerate}[start=3]
\item Podemos tomar $U_p$ abierto con clausura compacta dentro de una carta. Como $M$ es segundo numerable, existe un subcubrimiento numerable $\{U_{p_k}\}_{k\in\mathbb N}$.
\item Entonces
\end{enumerate}
\[
N=F(M)=\bigcup_{k\in\mathbb N}F(U_{p_k}).
\]
\begin{enumerate}[start=5]
\item Pero cada $F(U_{p_k})$ esta contenido, en coordenadas locales, en un subconjunto de la forma
\end{enumerate}
\[
\{(x_1,\ldots,x_r,0,\ldots,0)\}\subseteq\mathbb R^n,
\]
que tiene interior vacio si $r<n$. Esto contradice el teorema de Baire. Luego $r=n$ y $F$ es submersion.

\begin{itemize}
\item (c)
\end{itemize}
\begin{enumerate}
\item Si $F$ es inyectiva y no fuera inmersion, entonces $r<m$. En cartas de rango constante,
\end{enumerate}
\[
\psi\circ F\circ\varphi^{-1}(x_1,\ldots,x_m)=(x_1,\ldots,x_r,0,\ldots,0).
\]
\begin{enumerate}[start=2]
\item Tomando puntos distintos que coinciden en las primeras $r$ coordenadas y difieren en alguna de las restantes, sus imagenes por $F$ coinciden. Esto contradice que $F$ sea inyectiva. Por lo tanto $r=m$ y $F$ es inmersion.
\item Si $F$ es biyectiva, por 1 es submersion y por 2 es inmersion. Luego $r=n$ y $r=m$, asi que $m=n$. Entonces $(dF)_p$ es un isomorfismo para todo $p\in M$.
\item Por el teorema de la funcion inversa, $F$ es un difeomorfismo local. Como ademas $F$ es biyectiva, $F^{-1}$ esta bien definida globalmente y es suave localmente. Por lo tanto $F$ es un difeomorfismo.
\end{enumerate}
\bookqed
\end{bookproof}
\section{Cubo abierto y rebanadas}
\begin{bookdefinition}{Cubo abierto}
Sea $\epsilon > 0$. Venimos denotando por $C_\epsilon^n(0)$ al cubo abierto centrado en $0$ de $\mathbb R^n$ de ancho $2\epsilon$:

\[
C_\epsilon^n(0)=\{(x_1,\ldots,x_n)\in\mathbb R^n: |x_i|<\epsilon\}=(-\epsilon,\epsilon)\times\cdots\times(-\epsilon,\epsilon)\qquad (n\text{ veces}).
\]
\end{bookdefinition}
\begin{bookdefinition}{Rebanada de $C_{\epsilon}^{n}(0)$ a una altura}
Sea $m<n$ y sean $c_{m+1},\ldots,c_n$ constantes en $\mathbb R$ con $|c_j|<\epsilon$ para todo $j$.

\[
R(c_{m+1},\ldots,c_n):=\{y\in C_\epsilon^n(0):y_{m+1}=c_{m+1},\ldots,y_n=c_n\}.
\]
Se le llama rebanada de $C_\epsilon^n(0)$ a la altura de $(c_{m+1},\ldots,c_n)$.

\end{bookdefinition}
\begin{bookdefinition}{Rebanada de una variedad}
Sea $M$ una variedad de dimension $n$ y $(U,\varphi=(x_1,\ldots,x_n))$ una carta suave cubica de $M$ centrada en $p$; esto es $\varphi(U)=C_\epsilon^n(0)$ con $\varphi(p)=0$ de $\mathbb R^n$.

Dado $m\in\mathbb N$, $m<n$, y constantes $c_{m+1},\ldots,c_n$ con $|c_j|<\epsilon$, se define la rebanada $R(c_{m+1},\ldots,c_n)$ de $(U,\varphi)$ como el conjunto:

\[
R(c_{m+1},\ldots,c_n)=\{q\in U:x_{m+1}(q)=c_{m+1},\ldots,x_n(q)=c_n\}.
\]
Otra forma de verlo es

\[
R(c_{m+1},\ldots,c_{n})=\varphi ^{-1}((-\epsilon,\epsilon)^{m} \times \{ c_{m+1},\ldots c_{n} \})
\]
para esta idea es que se usa una carta cubica

\end{bookdefinition}
\begin{bookremark}{Ejercicio}
Para todo $p\in M$ variedad, existe una carta cubica centrada en $p$.

\end{bookremark}
\begin{bookproposition}{Ejercicio}
Sea $M$ variedad y $(U,\varphi)$ carta. Toda rebanada $R=R(c_{m+1},\ldots,c_n)$ con la topologia relativa de $M$ tiene una estructura diferenciable tal que $(R,i)$ es subvariedad de $U$ y a su vez de $M$, de dimension $m$. (Supongo $m<n$ por como definimos rebanada)

\end{bookproposition}
\begin{bookproof}{}
\begin{enumerate}
\item Consideremos la funcion
\end{enumerate}
\[
\pi_m:\mathbb R^n\to\mathbb R^m
\]
\[
(y_1,\ldots,y_n)\mapsto(y_1,\ldots,y_m)
\]
que proyecta a $(y_1,\ldots,y_n)$ en sus primeras $m$ coordenadas. 2. Supongamos como carta de $R$:

\[
(R,\widetilde\varphi=\pi_m\circ\varphi|_R)
\]
y la estructura dif. sobre $R$ seria el atlas maximal generado por $\{(R,\widetilde\varphi)\}$. 3. Notar que $\widetilde\varphi(R)$ es el abierto en $\mathbb R^m$

\[
C_\epsilon^m(0)=\{(y_1,\ldots,y_m)\in\mathbb R^m:|y_i|<\epsilon\}.
\]
aca estamos asumiendo que $\varphi$ es carta cubica, por definicion de rebanada 4. Veamos que $\widetilde\varphi$ es homeo entre $R$ y $C_\epsilon^m(0)$. 5. Claramente $\widetilde\varphi$ es continua, pues $\varphi|_R$ es continua, por tener $R$ la topologia relativa y por ser $\varphi$ continua 6. Ademas $\widetilde{\varphi}$ tiene inversa continua

\[
\varphi^{-1}\circ j:C_\epsilon^m(0)\to R
\]
donde $j$ es la funcion

\[
j:C_\epsilon^m(0)\to C_\epsilon^n(0),\qquad (y_1,\ldots,y_m)\mapsto(y_1,\ldots,y_m,c_{m+1},\ldots,c_n).
\]
\begin{enumerate}[start=7]
\item Por ultimo, veamos que $i:R\to M$ es suave y una inmersion. Para esto comparemos con las cartas $\widetilde\varphi$ y $\varphi$:
\end{enumerate}
\[
\varphi\circ i\circ\widetilde\varphi^{-1}:C_\epsilon^m(0)\subseteq\mathbb R^m\to C_\epsilon^n(0)\subseteq\mathbb R^n
\]
\[
(y_1,\ldots,y_m)\mapsto(y_1,\ldots,y_m,c_{m+1},\ldots,c_n).
\]
por lo que $i$ es suave 8. Y de paso resulta inmersion, pues $F=\varphi\circ i\circ\widetilde\varphi^{-1}$ es una inmersion.

\[
J(\varphi\circ i\circ\widetilde\varphi^{-1})=\begin{pmatrix}I_{m\times m}\\0\end{pmatrix}.
\]
por lo que $\operatorname{Imagen} d(\varphi\circ i\circ\widetilde\varphi^{-1})_{(y_1,\ldots,y_m)}$ tiene rango $m$, por lo que tiene kernel trivial (Teorema del rango y el kernel). Osea $(dF)_{(y_{1},\ldots,y_{m})}$ es inyectiva 9. Ademas esto implica que $(di)_q$ para todo $q\in R$ tambien es inyectiva, veamoslo 10. Notemos que $(d\varphi)_q\text{ y }(d\widetilde\varphi^{-1})_{\widetilde\varphi(q)}$ son isomorfismos por que $\varphi$ y $\widetilde{\varphi}$ son cartas osea difeomorfismos. 11. Entonces sea $v$ tal que $di(v)=0$ como $d \tilde{\varphi}^{-1}$ es isomorfismo existe $w$ tal que $d \tilde{\varphi}^{-1}(w)=v$ entonces

\[
(dF)_{q}(w)=d\varphi(di(d \tilde{\varphi}(w)))=d\varphi(di(w))=d\varphi(0)=0
\]
por que $d\varphi$ es es lineal y como $dF$ es iso entonces $w=0$ 12. Luego $0=d \tilde{\varphi}(w)=v$ mostrando que $\ker\{ (di)_{q} \}=0$ osea que $(di)_{q}$ es inyectiva, como queriamos

\bookqed
\end{bookproof}
